\section{持久化格式（FCS2）}

\subsection{从 FCSB 到 FCS2}
第~16~轮完成 \textbf{P2 格式 v2}（设计稿 \texttt{P2\_DESIGN.md}）：
\begin{itemize}[leftmargin=1.6em]
  \item 超级块魔数 \texttt{'FCS2'}，头扩至约 $48\,\mathrm{B}$ 量级字段布局（实现以头文件为准）。
  \item 清单项 $48\,\mathrm{B}$：\texttt{seq}/\texttt{ver} 升为 \textbf{64~位}，消除坑~102/103 类回绕静默错误的格式根因。
  \item 老库挂载可只读打开 v1，写路径触发升级；golden 镜像分 \texttt{tests/golden}（v2）与 \texttt{golden/v1}。
\end{itemize}

\subsection{超级块（双槽）}
\begin{itemize}[leftmargin=1.6em]
  \item 两槽交替；CRC 覆盖头 + 内联增量。
  \item 提交点语义：有效槽的状态字单次对齐写；挂载择优。
  \item ROUND17 起区分代数：\texttt{sb\_built\_}（锁内抓载荷）与 \texttt{sb\_synced\_}（锁外已持久）；页复用闸比较后者。
\end{itemize}

\subsection{清单操作码（节选）}
\begin{center}
\begin{tabular}{@{}clp{8.2cm}@{}}
\toprule
码 & 名 & 含义 \\
\midrule
1--5 & ADD/DEL/\ldots & run / split / seg 增减 \\
6 & \texttt{M\_RANGE} & 范围墓碑（只活在清单，不占键空间） \\
7 & \texttt{M\_RUN\_TRIM} & P2b：源 run 前缀释放；第~0~页为身份锚 \\
8 & \texttt{M\_RUN\_HOLE} & P2c 预留：run 内部释页（本冻结轮\textbf{未实现}） \\
\bottomrule
\end{tabular}
\end{center}

\subsection{页与记录}
\begin{itemize}[leftmargin=1.6em]
  \item 页大小 \texttt{PAGE}（与 XDB 共用布局常数）；数据页含 used、首末键、可选页内前缀（\texttt{pfx}/\texttt{fmt}）。
  \item 记录头瘦身：键 $8\,\mathrm{B}$ + 长度 $2\,\mathrm{B}$；版本是 run 属性；页 CRC 覆盖载荷。
  \item 大值：\texttt{len = OVFL\_MARK}，体为 $16\,\mathrm{B}$ 描述符，值在溢出链；\texttt{put} 时可先落链。
\end{itemize}

\subsection{身份与复用}
页号会复用，故段身份必须带 \texttt{seq}（及 rid）。提交失败前释放的页进 \texttt{pending\_free\_}，须待对应 SB 代数持久后才可复用——否则崩溃会破坏「完整旧态」。
